% Irodov Part VI §6.1 problems 1--18.
% Source: user-provided Arihant PDF pp.241--243 (printed pp.232--234),
% Answers pp.348--349, numerical input verified against page scans.
% Equations here use SI; C=Z1 Z2 e^2/(4 pi eps0) converts Gaussian e^2.
\atomicproblem{1}{232}
用汤姆孙的原子模型估算氢原子半径 $R$，
以及电子振动发出的光波长。已知电离能
$E_{\mathrm{ion}}=13.6\,\mathrm{eV}$。
\exsolution
汤姆孙模型将 $+e$ 均匀分布在半径 $R$ 的球内，
电子在球心稳定。电子由球心移至无穷远的
势能增加为
$E_{\mathrm{ion}}=3e^2/(8\pi\varepsilon_0R)$，
故
\[
 \boxed{R=\frac{3e^2}{8\pi\varepsilon_0E_{\mathrm{ion}}}
 \simeq0.159\,\mathrm{nm}}.
\]
在均匀正电球内电场 $E(r)=er/(4\pi\varepsilon_0R^3)$，
电子受到的回复力 $F=-e^2r/(4\pi\varepsilon_0R^3)$，
所以
\[
 \omega=\sqrt{\frac{e^2}{4\pi\varepsilon_0m_eR^3}},
 \quad\boxed{\lambda=\frac{2\pi c}{\omega}
 \simeq0.24\,\mu\mathrm m}.
\]
\excheck 本题严格依\emph{经典汤姆孙模型}计算，
不应拿现代量子氢原子的 Bohr 半径替代 $R$。

\atomicproblem{2}{232}
动能 $T=0.27\,\mathrm{MeV}$ 的 $\alpha$ 粒子
在金箔中被偏转 $\theta=60^\circ$。
求对应的碰撞瞄准参数 $b$。
\exsolution
对固定重核的库仑散射采用卢瑟福关系
\[
 \tan\frac\theta2=
 \frac{C}{2bT},\qquad
 C=\frac{(2e)(79e)}{4\pi\varepsilon_0}.
\]
故
\[
 \boxed{b=\frac C{2T}\cot\frac\theta2}.
\]
取 $e^2/(4\pi\varepsilon_0)
=1.440\,\mathrm{MeV\,fm}$ 得
\[
 \boxed{b\simeq7.30\times10^2\,\mathrm{fm}
 =0.730\,\mathrm{pm}}.
\]
\excheck 底本第六篇采用高斯单位制；
本译稿用 SI，故 $C$ 中明确包含 $1/(4\pi\varepsilon_0)$。

\atomicproblem{3}{232--233}
动能 $T=0.40\,\mathrm{MeV}$ 的 $\alpha$ 粒子
正面撞向 (a) 固定铅核；(b) 起初静止的\emph{自由}
$^{7}\mathrm{Li}$ 核。分别求最近距离。
\exsolution
设库仑能常数
$C=(2Z)e^2/(4\pi\varepsilon_0)$。
(a) 固定核没有反冲，初始动能在最近点全转为静电势能，
$T=C/r_{\min}$，其中 $Z_{\mathrm{Pb}}=82$。
因此
\[
 \boxed{r_{\min,\mathrm{Pb}}=\frac{2Z_{\mathrm{Pb}}k_e e^2}{T}
 \simeq0.590\,\mathrm{pm}}.
\]
(b) 自由锂核也会反冲，质心系相对动能为
$T_{\rm rel}=T\,M_{\rm Li}/(M_{\rm Li}+m_\alpha)$。
因此
\[
 \boxed{r_{\min,\mathrm{Li}}=
 \frac{2Z_{\mathrm{Li}}k_e e^2}{T}
 \left(1+\frac{m_\alpha}{M_{\mathrm{Li}}}\right)
 \simeq0.034\,\mathrm{pm}},
 \quad Z_{\mathrm{Li}}=3.
\]
\excheck $k_e=1/(4\pi\varepsilon_0)$；
自由靶核的修正因子大于1，不能直接用静止不动核的公式。

\atomicproblem{4}{233}
动能 $T=0.50\,\mathrm{MeV}$ 的 $\alpha$ 粒子
在固定汞核（$Z=80$）的库仑场中偏转
$\theta=90^\circ$。求轨道最小曲率半径和最近距离。
\exsolution
仍记 $C=2Zk_e e^2$。由库仑双曲线轨道几何
及卢瑟福散射公式可得
\[
 \boxed{\mathcal R_{\min}=
 \frac C{2T}\cot^2\frac\theta2},\qquad
 \boxed{r_{\min}=
 \frac C{2T}\left(1+\csc\frac\theta2\right)}.
\]
这里 $C=230.4\,\mathrm{MeV\,fm}$，
$\cot45^\circ=1$，
\[
 \boxed{\mathcal R_{\min}\simeq0.230\,\mathrm{pm}},
 \qquad\boxed{r_{\min}\simeq0.556\,\mathrm{pm}}.
\]
\excheck 两种长度不相同；
最近曲率半径属于轨迹局部几何，最近距离是到核的径向距离。

\atomicproblem{5}{233}
动能为 $T$、瞄准参数为 $b$ 的质子
被起初静止的金核散射，求传给金核的动量大小。
采用重核近似。
\exsolution
入射质子动量为 $p=\sqrt{2m_pT}$，
散射前后速率在固定靶近似下不变，
动量变化量为
$|\Delta\boldsymbol p|=2p\sin(\theta/2)$。
由
$\tan(\theta/2)=C/(2bT)$、
$C=79k_e e^2$ 得
\[
 \sin\frac\theta2=
 \frac1{\sqrt{1+(2bT/C)^2}}.
\]
根据动量守恒，金核所获动量大小为
\[
 \boxed{P_{\rm Au}=
 \frac{2\sqrt{2m_pT}}{\sqrt{1+(2bT/C)^2}}},
 \qquad C=79k_e e^2.
\]
\excheck 核的反冲动量方向与质子动量变化反向；
若求精细反冲能量，需要采用两体约化质量修正。

\atomicproblem{6}{233}
动能 $T=10\,\mathrm{MeV}$ 的质子以直线轨迹
掠过静止自由电子，最近距离 $b=10\,\mathrm{pm}$。
假定质子路径不偏折、电子在通过过程中几乎不动。
求电子得到的能量。
\exsolution
令质子速度
$v=\sqrt{2T/m_p}$。
电子的横向力为
$F_\perp(t)=k_ee^2 b/(b^2+v^2t^2)^{3/2}$。
沿完整经过时间积分，
\[
 \Delta p_\perp=\int_{-\infty}^{+\infty}F_\perp(t)\,dt
 =\frac{2k_ee^2}{bv}.
\]
电子获得的动能
\[
 T_e=\frac{(\Delta p_\perp)^2}{2m_e}
 =\boxed{\frac{m_p(k_ee^2)^2}{m_eb^2T}
 \simeq4.0\,\mathrm{eV}}.
\]
\excheck 纵向冲量为奇函数积分，等于0；
此题是小反冲、直线冲量近似而非完整双体散射。

\atomicproblem{7}{233}
粒子动能为 $T$、瞄准参数为 $b$，
穿过半径 $R$、深度 $U_0$ 的球形方势阱：
$U=0$（$r>R$）、$U=-U_0$（$r<R$）。
求偏转角 $\theta$ 与 $b$ 的关系。
\exsolution
粒子在阱内的速度比阱外大，
利用轨道的切向动量（角动量）守恒，
等效折射率
$n=v_{\rm in}/v_{\rm out}
=\sqrt{1+U_0/T}>1$。
令入射线与界面法线的夹角为 $i$，
内部为 $r$，有
$\sin i=b/R$、$\sin r=b/(nR)$。
两次穿越界面使轨迹总偏折
\[
 \boxed{\theta=
 2\left[\arcsin\frac bR-
 \arcsin\frac b{nR}\right]},
 \qquad 0\le b\le R.
\]
反解
\[
 \boxed{b=\frac{nR\sin(\theta/2)}
 {\sqrt{1+n^2-2n\cos(\theta/2)}}}.
\]
若 $b>R$ 则不进入阱，偏转角为0。
\excheck $U_0\to0$ 时 $n\to1$，
$\theta\to0$，符合无势阱极限。

\atomicproblem{8}{233}
半径 $R$ 的固定硬球受到半径为 $r$ 的
均匀平行粒子流照射，粒子与球作理想弹性碰撞。
求：(a) 偏角 $\theta$ 与瞄准参数 $b$；
(b) 碰后粒子落在 $(\theta,\theta+d\theta)$ 的条件概率；
(c) 偏角小于 $\pi/2$ 的概率。
\exsolution
把有限尺寸的入射粒子与球等效为
有效半径 $a=R+r$ 的点粒子镜面碰撞。
几何法线使碰前碰后速度对称，故
\[
 \boxed{\cos\frac\theta2=\frac ba},\quad 0\le b\le a.
\]
对于均匀入射、已发生碰撞的粒子，
撞击参数分布是
$dP=2b\,db/a^2$。
由 $b=a\cos(\theta/2)$ 微分并取绝对值，
\[
 \boxed{dP=\frac12\sin\theta\,d\theta}.
\]
积分 $0\le\theta<\pi/2$，
\[
 \boxed{P(\theta<\pi/2)=
 \int_0^{\pi/2}\tfrac12\sin\theta\,d\theta=\frac12}.
\]
\excheck (b)(c) 的概率以\emph{发生碰撞}为条件，
不涉及整个无限截面入射束中未碰撞的粒子。

\atomicproblem{9}{233}
动能 $T=1.0\,\mathrm{MeV}$ 的窄 $\alpha$ 束
正入射厚 $d=1.0\,\mu\mathrm m$ 的铂箔。
一个探测计数器位于偏转角 $60^\circ$、距靶
$\ell=10\,\mathrm{cm}$ 处，入口面积
$S=1.0\,\mathrm{cm^2}$。
求总入射粒子中进入计数器的比例。
\exsolution
铂的原子序数 $Z=78$、
密度 $\rho=21.5\,\mathrm{g/cm^3}$、
摩尔质量 $M\simeq195\,\mathrm{g/mol}$。
单位面积靶核数
$n_A=\rho dN_A/M$。
在单次散射、薄靶和小探测立体角近似下
\[
 \frac{\Delta N}{N}=
 n_A\left(\frac C{4T}\right)^2
 \csc^4\frac\theta2\,\Delta\Omega,\qquad
 C=2Zk_ee^2,\quad \Delta\Omega\simeq\frac S{\ell^2}.
\]
代 $\Delta\Omega=0.010\,\mathrm{sr}$ 得
\[
 \boxed{\frac{\Delta N}{N}\simeq3.3\times10^{-5}}.
\]
\excheck 必须将面积、长度和靶材面核密度统一单位；
此为\emph{探测几何}内的概率而不是总散射概率。

\atomicproblem{10}{233--234}
动能 $T=0.50\,\mathrm{MeV}$、通量
$\dot N_0=5.0\times10^3\,\mathrm{s^{-1}}$
的窄 $\alpha$ 束正入射金箔。
在 $\theta=60^\circ$、距靶 $\ell=15\,\mathrm{cm}$
处，测得粒子流密度
$J=40\,\mathrm{cm^{-2}s^{-1}}$。
求金箔厚度 $d$。
\exsolution
靶核体积数密度
$n_V=\rho_{\rm Au}N_A/M_{\rm Au}$，
单次卢瑟福散射预测的计数器位置通量为
\[
 J=\frac{\dot N_0}{\ell^2}
 n_Vd\left(\frac C{4T}\right)^2
 \csc^4\frac\theta2,\quad C=2(79)k_ee^2.
\]
故
\[
 \boxed{d=\frac{J\ell^2}{\dot N_0n_V}
 \left(\frac{4T}{C}\right)^2
 \sin^4\frac\theta2}.
\]
直接代入\emph{本版印刷题面}
$\rho_{\rm Au}\simeq19.3\,\mathrm{g/cm^3}$ 得
\[
 \boxed{d\simeq1.47\times10^{-4}\,\mathrm m
 =147\,\mu\mathrm m}.
\]
\excheck 本版书末印\emph{约 $1.5\,\mu\mathrm m$}，
与题面按上述薄靶公式计算相差\emph{约100倍}；
题面 $\dot N_0$ 与 $J$ 已按扫描页复查，暂保留该差异。
由此反推的厚箔已可能违反“单次散射薄靶”前提，
因此该数值不得标记为完全物理自洽。

\atomicproblem{11}{234}
一束 $\alpha$ 粒子依次照射质量面密度相同的
银箔和铂箔，后者相同角区计数率为前者的
$\eta=1.52$ 倍。
求铂原子序数（已知银原子序数与两材料原子量）。
\exsolution
相同质量面密度时核的面数密度与摩尔质量
$M$ 成反比；同样束流、能量与角度下
卢瑟福散射微分截面与 $Z^2$ 成正比。
因而
\[
 \eta=\frac{Z_{\rm Pt}^2/M_{\rm Pt}}
 {Z_{\rm Ag}^2/M_{\rm Ag}},\quad
 \boxed{Z_{\rm Pt}=Z_{\rm Ag}
 \sqrt{\eta\,\frac{M_{\rm Pt}}{M_{\rm Ag}}}}.
\]
以 $Z_{\rm Ag}=47$、
$M_{\rm Pt}\simeq195$、
$M_{\rm Ag}\simeq108$ 代入，
\[
 \boxed{Z_{\rm Pt}\simeq78}.
\]
\excheck 本版原题为 $\eta=1.52$，
不是 PDF 文本层中容易误读的“152”。

\atomicproblem{12}{234}
$\alpha$ 粒子束动能
$T=0.50\,\mathrm{MeV}$、
入射率 $I_0=5.0\times10^6\,\mathrm{s^{-1}}$，
正入射质量面密度
$\rho d=1.5\,\mathrm{mg/cm^2}$ 的金箔，
持续 $\tau=30\,\mathrm{min}$。
求散射角落在(a) $59^\circ$--$61^\circ$；
(b) 大于 $60^\circ$ 的粒子数。
\exsolution
靶核面数密度为
$n_A=\rho dN_A/M_{\rm Au}$，
总入射数 $N_0=I_0\tau$。
卢瑟福微分截面
$d\sigma/d\Omega=(C/(4T))^2\csc^4(\theta/2)$，
$C=2(79)k_ee^2$。
在极角区间 $\theta_a\le\theta\le\theta_b$
对全方位积分，得
\[
 \sigma_{ab}=\pi\left(\frac C{2T}\right)^2
 \left[\cot^2\frac{\theta_a}{2}
 -\cot^2\frac{\theta_b}{2}\right].
\]
因此
\[
 \boxed{N_{ab}=I_0\tau n_A\sigma_{ab}}.
\]
按\emph{扫描原题数值}代入，
\[
 \boxed{N_{59^\circ\text{--}61^\circ}
 \simeq1.62\times10^7},\qquad
 \boxed{N_{\theta>60^\circ}\simeq2.01\times10^8}.
\]
\excheck 本版书末分别列约 $1.6\times10^6$、
$2.0\times10^7$，低于照印刷题面算出的结果10倍。
原书扫描已确认质量面密度为 $1.5\,\mathrm{mg/cm^2}$，
故暂按原题算式并登记书末冲突，不自行把参数改小10倍。

\atomicproblem{13}{234}
速度 $v=6\times10^6\,\mathrm{m/s}$ 的窄质子束，
正入射厚 $d=1.0\,\mu\mathrm m$ 的银箔。
求后向半球（散射角 $\theta>90^\circ$）
的散射概率。
\exsolution
质子动能
$T=\tfrac12m_pv^2\simeq0.188\,\mathrm{MeV}$。
取银 $Z=47$，
$C=Zk_ee^2$，
靶核面数密度
$n_A=\rho_{\rm Ag}dN_A/M_{\rm Ag}$。
对 $\theta>90^\circ$ 积分卢瑟福截面：
\[
 \sigma_{\mathrm{back}}=
 \pi\left(\frac C{2T}\right)^2\cot^2(45^\circ)
 =\pi\left(\frac C{2T}\right)^2.
\]
故
\[
 \boxed{P=n_A\sigma_{\mathrm{back}}
 \simeq6.0\times10^{-3}}.
\]
\excheck 概率小于1，单次散射近似在此角区合理。

\atomicproblem{14}{234}
动能 $T=600\,\mathrm{keV}$ 的 $\alpha$ 束
正入射金箔，面核密度为
$n_A=1.1\times10^{19}\,\mathrm{cm^{-2}}$。
求偏转角小于 $\theta_0=20^\circ$
的束流份额（包括几乎不被散射的直行粒子）。
\exsolution
超过阈角的单次散射截面为
\[
 \sigma_{>\theta_0}=
 \pi\left(\frac C{2T}\right)^2
 \cot^2\frac{\theta_0}{2},\quad C=2(79)k_ee^2.
\]
薄靶近似下偏角超过阈值的粒子份额
$P_>=n_A\sigma_{>\theta_0}$。
因此偏角\emph{小于}阈值的份额为
\[
 \boxed{P_<\simeq1-n_A\pi\left(\frac C{2T}\right)^2
 \cot^2\frac{\theta_0}{2}\simeq0.60}.
\]
\excheck 单纯从 $\theta=0$ 积分卢瑟福公式会发散；
应统计“没有发生大角散射”的补事件。

\atomicproblem{15}{234}
动能 $T=1.4\,\mathrm{MeV}$ 的质子束照射黄铜箔，
质量面密度 $\rho d=1.5\,\mathrm{mg/cm^2}$，
铜与锌质量比为 $7:3$。
求偏角超过 $\theta_0=30^\circ$ 的粒子份额。
\exsolution
对元素 $j=\mathrm{Cu},\mathrm{Zn}$，阈角以上截面
\[
 \sigma_j(>\theta_0)=
 \pi\left(\frac{Z_jk_ee^2}{2T}\right)^2
 \cot^2\frac{\theta_0}{2}.
\]
两成分单位面积核数
$n_{A,j}=w_j(\rho d)N_A/M_j$，
$w_{\rm Cu}=0.7,w_{\rm Zn}=0.3$，所以
\[
 \boxed{P_>=(\rho d)N_A
 \frac{\pi(k_ee^2)^2}{4T^2}\cot^2\frac{\theta_0}{2}
 \left(\frac{0.7Z_{\rm Cu}^2}{M_{\rm Cu}}+
 \frac{0.3Z_{\rm Zn}^2}{M_{\rm Zn}}\right)}.
\]
用 $Z_{\rm Cu}=29,Z_{\rm Zn}=30$、
$M_{\rm Cu}\simeq63.55\,\mathrm{g/mol}$、
$M_{\rm Zn}\simeq65.38\,\mathrm{g/mol}$，
\[
 \boxed{P_>\simeq1.4\times10^{-3}}.
\]
\excheck 黄铜两种元素应按\emph{质量份额/摩尔质量}
分别计算面核密度，不能直接按0.7、0.3平均原子序数。

\atomicproblem{16}{234}
求 $\alpha$ 粒子动能 $T=1.5\,\mathrm{MeV}$
在铀核（$Z=92$）库仑场散射偏角超过
$\theta_0=60^\circ$ 的有效截面。
\exsolution
积分卢瑟福微分截面，
\[
 \boxed{\sigma_{>\theta_0}=
 \pi\left(\frac{2Zk_ee^2}{2T}\right)^2
 \cot^2\frac{\theta_0}{2}}.
\]
由于 $k_ee^2=1.440\,\mathrm{MeV\,fm}$，
代入 $Z=92$、$\theta_0=60^\circ$，
得到
\[
 \boxed{\sigma_{>60^\circ}\simeq
 7.35\times10^{-26}\,\mathrm{m^2}
 =0.735\,\mathrm{kb}}.
\]
\excheck $1\,\mathrm{kb}=10^{-25}\,\mathrm{m^2}$；
本题需作角积分而非只报某个散射角的微分截面。

\atomicproblem{17}{234}
实验测得单个金核将单色 $\alpha$ 粒子散射在
$90^\circ$--$180^\circ$ 间的有效截面为
$\sigma=0.50\,\mathrm{kb}$。
求：(a) $\alpha$ 的动能；
(b) $\theta=60^\circ$ 时微分截面
$d\sigma/d\Omega$（$\mathrm{kb/sr}$）。
\exsolution
记 $C=2(79)k_ee^2$，
在后向半球
\[
 \sigma_{>90^\circ}=
 \pi\left(\frac C{2T}\right)^2.
\]
由此
\[
 \boxed{T=\frac C2\sqrt{\frac{\pi}{\sigma}}
 \simeq0.90\,\mathrm{MeV}}.
\]
而在 $\theta=60^\circ$，
\[
 \frac{d\sigma}{d\Omega}
 =\left(\frac C{4T}\right)^2
 \csc^4(30^\circ)
 =\frac{4\sigma_{>90^\circ}}{\pi}.
\]
所以
\[
 \boxed{\left.\frac{d\sigma}{d\Omega}\right|_{60^\circ}
 \simeq0.637\,\mathrm{kb/sr}}.
\]
\excheck 角积分量 $\sigma$ 与微分量
$d\sigma/d\Omega$ 的物理单位不同，不能混用。

\atomicproblem{18}{234}
由经典辐射功率公式估计：振动角频率
$\omega=5.0\times10^{15}\,\mathrm{s^{-1}}$
的电子近似简谐振动，其机械能减少为初值的
$1/\eta$（$\eta=10$）需多久？
\exsolution
电子加速度为
$a(t)=-\omega^2x_0\cos\omega t$。
SI 制拉莫尔辐射功率
$P(t)=e^2a^2/(6\pi\varepsilon_0c^3)$。
对一个振动周期取平均，
$\langle a^2\rangle=\omega^4x_0^2/2$；
振动机械能 $E=\tfrac12m_e\omega^2x_0^2$。
所以
\[
 \frac{dE}{dt}=-\frac{e^2\omega^2}
 {6\pi\varepsilon_0m_ec^3}E
 =-\Gamma E,\qquad
 E(t)=E(0)e^{-\Gamma t}.
\]
于是
\[
 \boxed{t=\frac{6\pi\varepsilon_0m_ec^3}
 {e^2\omega^2}\ln\eta
 \simeq1.5\times10^{-8}\,\mathrm s=15\,\mathrm{ns}}.
\]
\excheck 这是小损耗、准简谐振动的周期平均模型，
并已将原书高斯单位的辐射公式转换为 SI。
